% Results Table - Global-APT Dataset
\begin{table}[htbp]
\centering
\caption{Performance comparison of anomaly detection algorithms on Global-APT dataset}
\label{tab:results}
\small
\begin{tabular}{lcccc}
\toprule
\textbf{Algorithm} & \textbf{ROC-AUC} & \textbf{Precision} & \textbf{Recall} & \textbf{F1-Score} \\
\midrule
GraphSAGE & \textbf{0.970} & 0.762 & \textbf{0.993} & \textbf{0.862} \\
NetWalk & 0.662 & 0.356 & 0.993 & 0.524 \\
SLADE & 0.556 & 0.344 & 0.996 & 0.511 \\
GAT & 0.521 & 0.337 & 1.000 & 0.504 \\
GeneralDYG & 0.494 & 1.000 & 0.004 & 0.008 \\
\bottomrule
\end{tabular}
\end{table}

% Evaluation Pipeline Description
The evaluation pipeline follows a standard workflow for temporal graph anomaly detection. We start with the raw JSON dataset containing 4,632 labeled events. The data is first converted to Parquet format for efficient processing, and IP addresses are mapped to numeric node IDs to enable graph operations. A chronological split is then performed, allocating 70\% of events to training, 15\% to validation, and 15\% to testing, preserving the temporal order of events. Each algorithm is trained on the training set, with hyperparameters tuned using the validation set. Final evaluation is performed on the test set, where we compute ROC-AUC, precision, recall, and F1-score. An important step in our pipeline is automatic threshold optimization, as we found that default thresholds (0.5) lead to poor performance, with most algorithms predicting all events as normal. We optimize thresholds by maximizing F1-score on the validation set, which significantly improves results across all algorithms.

% Commentary on results
The results in Table~\ref{tab:results} show that GraphSAGE achieves the best overall performance with a ROC-AUC of 0.97 and an F1-score of 0.86. Notably, it maintains high recall (99.3\%) while achieving reasonable precision (76.2\%), which is particularly challenging on imbalanced datasets where anomalies represent only 33\% of events.

The remaining algorithms (SLADE, GAT, NetWalk) exhibit a similar behavior pattern: they achieve near-perfect recall (99-100\%) but generate a substantial number of false positives, resulting in low precision (33-36\%). This trade-off prioritizes comprehensive anomaly detection over precision, which may be acceptable in security applications where missing an anomaly is more costly than investigating false alarms. However, the high false positive rate could limit practical deployment without additional filtering mechanisms.

GeneralDYG performs poorly on this dataset, predicting almost all events as normal. This suggests the model requires significant hyperparameter tuning or may not be well-suited for small, highly centralized graph structures. The model's extreme conservatism (only 1 positive prediction out of 696 test events) indicates a fundamental mismatch between its learning strategy and the dataset characteristics.

A key finding from our experiments is the critical importance of threshold optimization. Using default classification thresholds (0.5) resulted in most algorithms predicting all events as normal, rendering them effectively useless. After implementing automatic threshold optimization based on F1-score maximization, we observed dramatic improvements. For instance, GraphSAGE's optimal threshold is 0.61, while GAT requires a much lower threshold of 0.01 to achieve 100\% recall. This threshold optimization step was essential for obtaining meaningful results from all algorithms.
